% !TeX root = ../main.tex

\chapter{Darstellung ganzer Zahlen}

\section{$b$-adische Darstellung natürlicher Zahlen}

Natürliche Zahlen werden in der \textbf{Binärdarstellung}
(auch $2$-adische Darstellung genannt) gespeichert.
Diese ist analog zur gewohnten Dezimaldarstellung.

\ex{Dezimal- und Binärdarstellung}{
  Die Zahl $106$ lässt sich darstellen als
  \begin{align*}
    106
    &=1\cdot10^2+0\cdot10^1+6\cdot10^0\\
    &=1\cdot2^6+1\cdot2^5+0\cdot2^4
      +1\cdot2^3+0\cdot2^2+1\cdot2^1+0\cdot2^0.
  \end{align*}

  In Binärdarstellung schreibt man daher
  \[
    106=(1101010)_2.
  \]
}

Statt $10$ oder $2$ kann man jede natürliche Zahl
$b\ge2$ als Basis wählen.

\thm{$b$-adische Darstellung}{
  Seien $b,n\in\bbN$ mit $b\ge2$.

  Dann existieren eindeutig bestimmte Zahlen
  $l\in\bbN$ und
  \[
    z_i\in\{0,\dots,b-1\}
    \qquad (i=0,\dots,l-1)
  \]
  mit $z_{l-1}\ne0$, sodass
  \[
    \boxed{n=\sum_{i=0}^{l-1}z_i b^i.}
  \]

  Das Wort $z_{l-1}\cdots z_0$ heißt die
  \textbf{$b$-adische Darstellung} von $n$.

  Man schreibt auch
  \[
    n=(z_{l-1}\cdots z_0)_b.
  \]

  Für die Anzahl der Ziffern gilt
  \[
    \boxed{l=1+\lfloor\log_b n\rfloor.}
  \]
}

\nt{
  \textbf{Vollständige Induktion:}

  Um eine Aussage $A(n)$ für alle $n\in\bbN$
  zu beweisen, genügt es zu zeigen:
  \begin{enumerate}
    \item \textbf{Induktionsanfang:} $A(1)$ gilt.
    \item \textbf{Induktionsschritt:}
      Für alle $i\in\bbN$ folgt aus $A(i)$
      die Aussage $A(i+1)$.
  \end{enumerate}

  Beim Induktionsschritt nennt man $A(i)$
  die \textbf{Induktionsvoraussetzung}.

  Allgemeiner kann man für eine Menge $M$ eine
  Aussage $A(m)$ für alle $m\in M$ mittels einer
  Funktion $f\colon M\to\bbN$ beweisen:

  Man zeigt, dass aus der Annahme, $A(m)$ sei falsch,
  ein $m'\in M$ mit
  \[
    f(m')<f(m)
    \qquad\text{und}\qquad
    A(m')\text{ falsch}
  \]
  folgt.

  Dies bezeichnet man als \textbf{Induktion über $f$}.
}

\begin{myproof}
  Wir zeigen zunächst die Formel für die Anzahl der
  Ziffern, anschließend die Eindeutigkeit und Existenz.

  \medskip
  \textbf{1. Anzahl der Ziffern:}

  Sei
  \[
    n=\sum_{i=0}^{l-1}z_i b^i
  \]
  mit $z_i\in\{0,\dots,b-1\}$ und $z_{l-1}\ne0$.

  Dann gilt
  \begin{align*}
    b^{l-1}
    &\le n\\
    &\le\sum_{i=0}^{l-1}(b-1)b^i\\
    &=b^l-1.
  \end{align*}

  Somit folgt
  \[
    b^{l-1}\le n<b^l
  \]
  und daher
  \[
    l-1=\lfloor\log_b n\rfloor.
  \]

  \medskip
  \textbf{2. Eindeutigkeit:}

  Aus der Anzahl der Ziffern folgt, dass $l$
  eindeutig durch $n$ bestimmt ist.

  Für eine feste Länge $l$ ist die Abbildung
  \[
    (z_{l-1},\dots,z_0)
    \longmapsto
    \sum_{i=0}^{l-1}z_i b^i
  \]
  nach dem Lemma über die Kodierung von Wörtern
  injektiv.

  Daher sind auch die Ziffern eindeutig bestimmt.

  \medskip
  \textbf{3. Existenz:}

  Wir führen vollständige Induktion über
  \[
    l(n):=1+\lfloor\log_b n\rfloor.
  \]

  \textbf{Induktionsanfang:}

  Sei $l(n)=1$. Dann gilt
  \[
    n\in\{1,\dots,b-1\}.
  \]

  Mit $z_0:=n$ erhalten wir die Darstellung
  \[
    n=z_0b^0.
  \]

  \textbf{Induktionsvoraussetzung:}

  Für alle natürlichen Zahlen $m$ mit
  $l(m)<l(n)$ existiere eine $b$-adische
  Darstellung.

  \textbf{Induktionsschritt:}

  Sei nun $n\in\bbN$ mit $l(n)\ge2$.

  Setze
  \[
    n':=\left\lfloor\frac nb\right\rfloor.
  \]

  Dann gilt
  \[
    l':=l(n')=l(n)-1.
  \]

  Nach Induktionsvoraussetzung besitzt $n'$
  eine Darstellung
  \[
    n'=\sum_{i=0}^{l'-1}z'_i b^i
  \]
  mit
  \[
    z'_i\in\{0,\dots,b-1\},
    \qquad z'_{l'-1}\ne0.
  \]

  Wir definieren
  \[
    z_0:=n\bmod b
  \]
  und
  \[
    z_i:=z'_{i-1}
    \qquad(i=1,\dots,l').
  \]

  Damit folgt
  \begin{align*}
    n
    &=b\left\lfloor\frac nb\right\rfloor
      +(n\bmod b)\\
    &=bn'+z_0\\
    &=b\sum_{i=0}^{l'-1}z'_i b^i+z_0\\
    &=\sum_{i=1}^{l'}z'_{i-1}b^i+z_0\\
    &=\sum_{i=0}^{l-1}z_i b^i.
  \end{align*}

  Da $z_{l-1}=z'_{l'-1}\ne0$, ist dies eine
  gültige $b$-adische Darstellung.

  Damit sind Existenz und Eindeutigkeit bewiesen.
\end{myproof}


\subsection{Horner-Schema}

Liegt eine Zahl in $b$-adischer Darstellung
$(z_{l-1}\cdots z_0)_b$ vor, so kann man ihren
Wert mit dem \textbf{Horner-Schema} berechnen:

\[
  \begin{aligned}
    \sum_{i=0}^{l-1}z_i b^i
    =z_0+b\bigl(z_1+b\bigl(z_2+\cdots
    +b(z_{l-2}+bz_{l-1})\cdots\bigr)\bigr).
  \end{aligned}
\]

Dadurch lassen sich Multiplikationen einsparen.

\nt{
  Umgekehrt liefert der Existenzbeweis einen
  Algorithmus zur Berechnung der $b$-adischen
  Darstellung einer natürlichen Zahl:

  Man bestimmt wiederholt den Rest
  \[
    z_0=n\bmod b
  \]
  und ersetzt anschließend
  \[
    n\leftarrow\left\lfloor\frac nb\right\rfloor.
  \]

  Die aufeinanderfolgenden Reste ergeben die
  Ziffern der Darstellung in umgekehrter Reihenfolge.
}

\ex{Übliche Zahlensysteme}{
  Neben der Dezimaldarstellung werden insbesondere
  folgende Zahlensysteme verwendet:
  \begin{itemize}
    \item \textbf{Binärdarstellung:} $b=2$
    \item \textbf{Oktaldarstellung:} $b=8$
    \item \textbf{Hexadezimaldarstellung:} $b=16$
  \end{itemize}

  Für Basen $b>10$ verwendet man die Buchstaben
  $\mathrm{A},\mathrm{B},\mathrm{C},\dots$
  für Ziffern größer als $9$.
}
